% GENERATED FILE. Edit canonical lessons, then run npm run generate:books.
% canonical-source-hash: fnv1a64:feddda72be31b9c8
% canonical-lessons: JA-C203-hikouki, JA-C203-kuukou, JA-C203-tobu, JA-C203-zaseki, JA-C203-yuzuru, JA-R203-review-198-203

\chapter{Planes, Airports and Seats}
\label{ch:ja-planes-airports-and-seats}
\hlchaptermodality{203}

\begin{chapteropening}
\textbf{By the end of this chapter:} \emph{I can talk about flying and giving up a seat.}

Complete the last lesson of chapter 203: I can talk about flying and giving up a seat.
\end{chapteropening}

\section[\texorpdfstring{\ja{ひこうき} (hikōki)}{ひこうき (hikōki)}]{\ja{ひこうき} (hikōki) — a plane; an aeroplane}

\label{lesson:JA-C203-hikouki}



\begin{quote}
Before the new one: say the Japanese for to halt, then the Japanese for a car park.
\end{quote}

\subsection*{You'll want to know: \ja{ひこうき}}
\textbf{\ja{ひこうき}} — \emph{hikōki} — \textquotedblleft{}a plane; an aeroplane\textquotedblright{}.

\begin{quote}
\textbf{\ja{ひこうきで} \ja{しまに} \ja{いきます。}} — \emph{hikōki de shima ni ikimasu} — \textquotedblleft{}I go to the island by plane.\textquotedblright{}
\end{quote}

\begin{cousinweb}
\textbf{\ja{ひこうき}} (\emph{hikōki}) is Sino-Japanese (kango): three Chinese characters read with their on'yomi, readings borrowed from Chinese. They are \emph{hi}, \textquotedblleft{}to fly\textquotedblright{} (Middle Chinese \emph{pj+j}), \emph{kō}, \textquotedblleft{}to go\textquotedblright{} (\emph{haeng}) and \emph{ki}, \textquotedblleft{}machine\textquotedblright{} (\emph{kj+j}) (KANJIDIC2; Middle Chinese in Baxter's transcription).
\end{cousinweb}

\begin{grammarlens}[title={What you've built}]
One more word for everyday life at A2.
\end{grammarlens}

\subsection*{Your turn}
\begin{itemize}
\raggedright
  \item \emph{Say it:} \emph{hikōki}
  \item \emph{Say it:} \emph{hikōki}, once more
  \item \emph{Say it:} \emph{hikōki de shima ni ikimasu} — I go to the island by plane.
  \item \emph{Recall:} say the Japanese for to halt, then the Japanese for a car park, then say \emph{hikōki} again
\end{itemize}

\begin{tcolorbox}[breakable,colback=teal!4,colframe=teal!35!black,title={Before you move on}]
What does \ja{ひこうき} mean? (\textquotedblleft{}A plane; an aeroplane\textquotedblright{}.) Say it once more.
\end{tcolorbox}
\section[\texorpdfstring{\ja{くうこう} (kūkō)}{くうこう (kūkō)}]{\ja{くうこう} (kūkō) — an airport}

\label{lesson:JA-C203-kuukou}



\begin{quote}
Before the new one: say the Japanese for a car park, then the Japanese for a plane.
\end{quote}

\subsection*{You'll want to know: \ja{くうこう}}
\textbf{\ja{くうこう}} — \emph{kūkō} — \textquotedblleft{}an airport\textquotedblright{}.

\begin{quote}
\textbf{\ja{くうこうまで} \ja{ちかてつで} \ja{いきます。}} — \emph{kūkō made chikatetsu de ikimasu} — \textquotedblleft{}I go to the airport by subway.\textquotedblright{}
\end{quote}

\begin{cousinweb}
\textbf{\ja{くうこう}} (\emph{kūkō}) is Sino-Japanese (kango): two Chinese characters read with their on'yomi, readings borrowed from Chinese. They are \emph{kū}, \textquotedblleft{}empty, sky\textquotedblright{} (Middle Chinese \emph{khuwng}) and \emph{kō}, \textquotedblleft{}harbour\textquotedblright{} (\emph{kaewngX}) (KANJIDIC2; Middle Chinese in Baxter's transcription).
\end{cousinweb}

\begin{grammarlens}[title={What you've built}]
One more word for everyday life at A2.
\end{grammarlens}

\subsection*{Your turn}
\begin{itemize}
\raggedright
  \item \emph{Say it:} \emph{kūkō}
  \item \emph{Say it:} \emph{kūkō}, once more
  \item \emph{Say it:} \emph{kūkō made chikatetsu de ikimasu} — I go to the airport by subway.
  \item \emph{Recall:} say the Japanese for a car park, then the Japanese for a plane, then say \emph{kūkō} again
\end{itemize}

\begin{tcolorbox}[breakable,colback=teal!4,colframe=teal!35!black,title={Before you move on}]
What does \ja{くうこう} mean? (\textquotedblleft{}An airport\textquotedblright{}.) Say it once more.
\end{tcolorbox}
\section[\texorpdfstring{\ja{とぶ} (tobu)}{とぶ (tobu)}]{\ja{とぶ} (tobu) — to fly through the air; to fly}

\label{lesson:JA-C203-tobu}



\begin{quote}
Before the new one: say the Japanese for a plane, then the Japanese for an airport.
\end{quote}

\subsection*{You'll want to know: \ja{とぶ}}
\textbf{\ja{とぶ}} — \emph{tobu} — \textquotedblleft{}to fly through the air; to fly\textquotedblright{}.

\begin{quote}
\textbf{\ja{ひこうきが} \ja{うみの} \ja{うえを} \ja{とびます。}} — \emph{hikōki ga umi no ue o tobimasu} — \textquotedblleft{}The plane flies over the sea.\textquotedblright{}
\end{quote}

Polite \textbf{\ja{とびます}}, \emph{tobimasu}; a godan verb.

\begin{cousinweb}
\textbf{\ja{とぶ}} (\emph{tobu}) is a native Japanese word (wago), inherited rather than borrowed from Chinese; its kanji was attached later as a kun'yomi, and its older origin is not settled (JMdict, KANJIDIC2).
\end{cousinweb}

\begin{grammarlens}[title={What you've built}]
One more word for everyday life at A2.
\end{grammarlens}

\subsection*{Your turn}
\begin{itemize}
\raggedright
  \item \emph{Say it:} \emph{tobu}
  \item \emph{Say it:} \emph{tobu}, once more
  \item \emph{Say it:} \emph{hikōki ga umi no ue o tobimasu} — The plane flies over the sea.
  \item \emph{Recall:} say the Japanese for a plane, then the Japanese for an airport, then say \emph{tobu} again
\end{itemize}

\begin{tcolorbox}[breakable,colback=teal!4,colframe=teal!35!black,title={Before you move on}]
What does \ja{とぶ} mean? (\textquotedblleft{}To fly through the air; to fly\textquotedblright{}.) Say it once more.
\end{tcolorbox}
\section[\texorpdfstring{\ja{ざせき} (zaseki)}{ざせき (zaseki)}]{\ja{ざせき} (zaseki) — a seat}

\label{lesson:JA-C203-zaseki}



\begin{quote}
Before the new one: say the Japanese for an airport, then the Japanese for to fly through the air.
\end{quote}

\subsection*{You'll want to know: \ja{ざせき}}
\textbf{\ja{ざせき}} — \emph{zaseki} — \textquotedblleft{}a seat\textquotedblright{}.

\begin{quote}
\textbf{\ja{ざせきは} \ja{まどの} \ja{ちかくです。}} — \emph{zaseki wa mado no chikaku desu} — \textquotedblleft{}The seat is near the window.\textquotedblright{}
\end{quote}

\begin{cousinweb}
\textbf{\ja{ざせき}} (\emph{zaseki}) is Sino-Japanese (kango): two Chinese characters read with their on'yomi, readings borrowed from Chinese. They are \emph{za}, \textquotedblleft{}to sit\textquotedblright{} (Middle Chinese \emph{dzwaH}) and \emph{seki}, \textquotedblleft{}mat, seat\textquotedblright{} (\emph{zjek}) (KANJIDIC2; Middle Chinese in Baxter's transcription).
\end{cousinweb}

\begin{grammarlens}[title={What you've built}]
One more word for everyday life at A2.
\end{grammarlens}

\subsection*{Your turn}
\begin{itemize}
\raggedright
  \item \emph{Say it:} \emph{zaseki}
  \item \emph{Say it:} \emph{zaseki}, once more
  \item \emph{Say it:} \emph{zaseki wa mado no chikaku desu} — The seat is near the window.
  \item \emph{Recall:} say the Japanese for an airport, then the Japanese for to fly through the air, then say \emph{zaseki} again
\end{itemize}

\begin{tcolorbox}[breakable,colback=teal!4,colframe=teal!35!black,title={Before you move on}]
What does \ja{ざせき} mean? (\textquotedblleft{}A seat\textquotedblright{}.) Say it once more.
\end{tcolorbox}
\section[\texorpdfstring{\ja{ゆずる} (yuzuru)}{ゆずる (yuzuru)}]{\ja{ゆずる} (yuzuru) — to give up (a seat); to hand over}

\label{lesson:JA-C203-yuzuru}



\begin{quote}
Before the new one: say the Japanese for to fly through the air, then the Japanese for a seat.
\end{quote}

\subsection*{You'll want to know: \ja{ゆずる}}
\textbf{\ja{ゆずる}} — \emph{yuzuru} — \textquotedblleft{}to give up (a seat); to hand over\textquotedblright{}.

\begin{quote}
\textbf{\ja{でんしゃで} \ja{ざせきを} \ja{ゆずります。}} — \emph{densha de zaseki o yuzurimasu} — \textquotedblleft{}I give up my seat on the train.\textquotedblright{}
\end{quote}

Polite \textbf{\ja{ゆずります}}, \emph{yuzurimasu}; a godan verb.

\begin{cousinweb}
\textbf{\ja{ゆずる}} (\emph{yuzuru}) is a native Japanese word (wago), inherited rather than borrowed from Chinese; its kanji was attached later as a kun'yomi, and its older origin is not settled (JMdict, KANJIDIC2).
\end{cousinweb}

\begin{grammarlens}[title={What you've built}]
One more word for everyday life at A2.
\end{grammarlens}

\subsection*{Your turn}
\begin{itemize}
\raggedright
  \item \emph{Say it:} \emph{yuzuru}
  \item \emph{Say it:} \emph{yuzuru}, once more
  \item \emph{Say it:} \emph{densha de zaseki o yuzurimasu} — I give up my seat on the train.
  \item \emph{Recall:} say the Japanese for to fly through the air, then the Japanese for a seat, then say \emph{yuzuru} again
\end{itemize}

\begin{tcolorbox}[breakable,colback=teal!4,colframe=teal!35!black,title={Before you move on}]
What does \ja{ゆずる} mean? (\textquotedblleft{}To give up (a seat); to hand over\textquotedblright{}.) Say it once more.
\end{tcolorbox}
\section[(dialogue)]{Trains, Roads and Planes - a review}

\label{lesson:JA-R203-review-198-203}



\begin{quote}
Say the Japanese for to ride and to get off.
\end{quote}

\subsection*{Your turn}
Say these aloud:

\begin{itemize}
\raggedright
  \item to ride, to get off, a bicycle, a subway, to get to
  \item departure, arrival, a limited express, to come to a stop, an express
  \item to change trains, to leave, a railway, a boarding point, a bus stop
  \item to cross, to be crowded, to cut across, to rush out, to bump into
  \item an automobile, driving, a driver, to halt, a car park
  \item a plane, an airport, to fly through the air, a seat, to give up
\end{itemize}

\begin{tcolorbox}[breakable,colback=teal!4,colframe=teal!35!black,title={Before you move on}]
To ride? (\textbf{\ja{のる}}, \emph{noru}.) To get off? (\textbf{\ja{おりる}}, \emph{oriru}.) A bicycle? (\textbf{\ja{じてんしゃ}}, \emph{jitensha}.) A subway? (\textbf{\ja{ちかてつ}}, \emph{chikatetsu}.) To get to? (\textbf{\ja{つく}}, \emph{tsuku}.) Departure? (\textbf{\ja{しゅっぱつ}}, \emph{shuppatsu}.) Arrival? (\textbf{\ja{とうちゃく}}, \emph{tōchaku}.) A limited express? (\textbf{\ja{とっきゅう}}, \emph{tokkyū}.) To come to a stop? (\textbf{\ja{とまる}}, \emph{tomaru}.) An express? (\textbf{\ja{きゅうこう}}, \emph{kyūkō}.) To change trains? (\textbf{\ja{のりかえる}}, \emph{norikaeru}.) To leave? (\textbf{\ja{でる}}, \emph{deru}.) A railway? (\textbf{\ja{てつどう}}, \emph{tetsudō}.) A boarding point? (\textbf{\ja{のりば}}, \emph{noriba}.) A bus stop? (\textbf{\ja{ていりゅうじょ}}, \emph{teiryūjo}.) To cross? (\textbf{\ja{わたる}}, \emph{wataru}.) To be crowded? (\textbf{\ja{こむ}}, \emph{komu}.) To cut across? (\textbf{\ja{よこぎる}}, \emph{yokogiru}.) To rush out? (\textbf{\ja{とびだす}}, \emph{tobidasu}.) To bump into? (\textbf{\ja{ぶつかる}}, \emph{butsukaru}.) An automobile? (\textbf{\ja{じどうしゃ}}, \emph{jidōsha}.) Driving? (\textbf{\ja{うんてん}}, \emph{unten}.) A driver? (\textbf{\ja{うんてんしゅ}}, \emph{untenshu}.) To halt? (\textbf{\ja{とめる}}, \emph{tomeru}.) A car park? (\textbf{\ja{ちゅうしゃじょう}}, \emph{chūshajō}.) A plane? (\textbf{\ja{ひこうき}}, \emph{hikōki}.) An airport? (\textbf{\ja{くうこう}}, \emph{kūkō}.) To fly through the air? (\textbf{\ja{とぶ}}, \emph{tobu}.) A seat? (\textbf{\ja{ざせき}}, \emph{zaseki}.) To give up? (\textbf{\ja{ゆずる}}, \emph{yuzuru}.)
\end{tcolorbox}
